\section*{Bab \ref{ch:g5:water-shapes-the-land} --- Air Membentuk Daratan}

\begin{solution}{exo:g5:water-shapes-the-land:1}
(a) \omterm{def:g5:water-shapes-the-land:weathering}{Pelapukan}: \omterm{def:g3:stones-pebbles-sand:rock}{batuan} pecah di tempatnya. (b) \omterm{def:g5:water-shapes-the-land:erosion}{Erosi}: \omterm{def:g3:stones-pebbles-sand:sand}{pasir} dikikis dan
diangkut pergi. (c) \omterm{def:g5:water-shapes-the-land:deposition}{Pengendapan}: \omterm{def:g3:stones-pebbles-sand:sand}{lumpur} mengendap. (d) \omterm{def:g5:water-shapes-the-land:weathering}{Pelapukan}: akar
memecahkan batu di tempatnya.
\end{solution}

\begin{solution}{exo:g5:water-shapes-the-land:2}
\omterm{def:g5:water-shapes-the-land:source-mouth}{Hulu} adalah tempat \omterm{def:g1:landscapes-around-us:water}{sungai} bermula (mata air, gletser yang mencair, daratan
tinggi yang basah). \omterm{def:g5:water-shapes-the-land:source-mouth}{Muara} adalah tempat \omterm{def:g1:landscapes-around-us:water}{sungai} berakhir, di \omterm{def:g2:water-on-earth:ocean}{laut}, \omterm{def:g1:landscapes-around-us:water}{danau},
atau \omterm{def:g1:landscapes-around-us:water}{sungai} yang lebih besar. \omterm{def:g5:water-shapes-the-land:source-mouth}{Anak sungai} adalah \omterm{def:g1:landscapes-around-us:water}{sungai} yang lebih kecil
yang mengalir masuk ke \omterm{def:g1:landscapes-around-us:water}{sungai} yang lebih besar.
\end{solution}

\begin{solution}{exo:g5:water-shapes-the-land:3}
\omterm{def:g5:water-shapes-the-land:source-mouth}{Hulu}; \omterm{def:g1:landscapes-around-us:landform}{lembah} sempit berbentuk V; \omterm{def:g5:water-shapes-the-land:meander}{meander}; \omterm{def:g5:water-shapes-the-land:delta}{delta}.
\end{solution}

\begin{solution}{exo:g5:water-shapes-the-land:4}
(a) Salah: embun beku, air, akar, panas, dan dingin perlahan-lahan
memecahkan \omterm{def:g3:stones-pebbles-sand:rock}{batuan} yang keras sekalipun (\omterm{def:g5:water-shapes-the-land:weathering}{pelapukan}). (b) Benar: makin deras
airnya, makin besar kepingan yang dapat dipindahkannya. (c) Salah: kepingan
yang paling berat mengendap lebih dahulu; \omterm{def:g3:stones-pebbles-sand:sand}{lumpur} yang paling halus
mengendap paling akhir, di tempat airnya hampir diam, di \omterm{def:g1:landscapes-around-us:water}{danau} atau di \omterm{def:g2:water-on-earth:ocean}{laut}.
\end{solution}

\begin{solution}{exo:g5:water-shapes-the-land:5}
A terletak di sisi luar tikungan: airnya di sana deras dan dalam, dan air
itu menggerus tepi \omterm{def:g1:landscapes-around-us:water}{sungai} (\omterm{def:g5:water-shapes-the-land:erosion}{erosi}). B terletak di sisi dalam tikungan: airnya
di sana lambat, dan air itu menjatuhkan \omterm{def:g3:stones-pebbles-sand:sand}{pasir} (\omterm{def:g5:water-shapes-the-land:deposition}{pengendapan}).
\end{solution}

\begin{solution}{exo:g5:water-shapes-the-land:6}
Di \omterm{def:g4:world-physical-map:range}{pegunungan} \omterm{def:g1:landscapes-around-us:water}{sungai} itu curam dan deras: \omterm{def:g1:landscapes-around-us:water}{sungai} dapat menggulingkan batu
besar, dan butiran kecil dibawanya pergi. Di dekat \omterm{def:g2:water-on-earth:ocean}{laut} \omterm{def:g1:landscapes-around-us:water}{sungai} itu lambat:
\omterm{def:g1:landscapes-around-us:water}{sungai} tidak lagi dapat memindahkan batu besar, hanya \omterm{def:g3:stones-pebbles-sand:sand}{pasir} dan \omterm{def:g3:stones-pebbles-sand:sand}{lumpur}, dan
\omterm{def:g1:landscapes-around-us:water}{sungai} menjatuhkannya di sana. Selain itu, dalam perjalanannya batu-batu
itu sudah pecah dan aus menjadi butiran yang lebih kecil.
\end{solution}

\begin{solution}{exo:g5:water-shapes-the-land:7}
(a) Aliran \omterm{def:g5:water-shapes-the-land:source-mouth}{hulu}: $600 \div 6 = 100$ kilometer. Aliran tengah:
$600 \div 2 = 300$ kilometer. (b) Aliran hilir: $600 - 100 - 300 = 200$
kilometer. Karena $3 \times 200 = 600$, aliran hilir adalah $\frac{1}{3}$
dari \omterm{def:g1:landscapes-around-us:water}{sungai} itu. (Periksa: $\frac{1}{6} + \frac{1}{2} + \frac{1}{3} = \frac{1}{6} +
\frac{3}{6} + \frac{2}{6} = \frac{6}{6}$.)
\end{solution}

\begin{solution}{exo:g5:water-shapes-the-land:8}
Bercak-bercak pucat itu adalah \omterm{def:g3:stones-pebbles-sand:sand}{pasir} dan \omterm{def:g3:stones-pebbles-sand:pebble}{kerikil} yang diendapkan di sisi
dalam tikungan, tempat airnya lambat (\omterm{def:g5:water-shapes-the-land:deposition}{pengendapan}). \omterm{def:g1:landscapes-around-us:water}{Danau-danau} gelap yang
melengkung itu adalah kelokan lama \omterm{def:g1:landscapes-around-us:water}{sungai}: dua kelokan tumbuh begitu dekat
sehingga \omterm{def:g1:landscapes-around-us:water}{sungai} menerobos leher daratan di antara keduanya, mengambil jalan
pintas, dan meninggalkan kelokan lama sebagai \omterm{def:g1:landscapes-around-us:water}{danau}.
\end{solution}

\begin{solution}{exo:g5:water-shapes-the-land:9}
Padang rumput datar di samping \omterm{def:g1:landscapes-around-us:water}{sungai} itu adalah \omterm{def:g5:water-shapes-the-land:flood}{dataran banjirnya}: daratan
yang digenangi \omterm{def:g1:landscapes-around-us:water}{sungai} ketika \omterm{def:g5:water-shapes-the-land:flood}{banjir}. Padang rumput itu datar justru karena
\omterm{def:g1:landscapes-around-us:water}{sungai} membangunnya dari \omterm{def:g3:stones-pebbles-sand:sand}{lumpurnya} sendiri. Suatu hari \omterm{def:g1:landscapes-around-us:water}{sungai} akan \omterm{def:g5:water-shapes-the-land:flood}{banjir}
lagi dan sekolah itu akan terisi air, dan air \omterm{def:g5:water-shapes-the-land:flood}{banjir} yang deras itu
berbahaya. Sekolah sebaiknya dibangun di tanah yang lebih tinggi, dan
padang rumput itu dibiarkan sebagai ladang, padang rumput, atau taman.
\end{solution}

\begin{solution}{exo:g5:water-shapes-the-land:10}
(a) Satu tahun ada 12 bulan: $12$ miliar $\div\, 12 = 1$ miliar ton setiap
bulan, rata-rata. (b) Satu miliar adalah $1\,000\,000\,000$; dibagi 10,
hasilnya $100\,000\,000$: seratus juta muatan truk.
\end{solution}

\begin{solution}{exo:g5:water-shapes-the-land:11}
Salah satu kemungkinan ceritanya. Butir itu adalah bagian dari \omterm{def:g3:stones-pebbles-sand:rock}{batuan}
tinggi di sebuah \omterm{def:g1:landscapes-around-us:landform}{gunung}. Embun beku membuka retakan di sekitarnya, malam
demi malam, sampai sekeping \omterm{def:g3:stones-pebbles-sand:rock}{batuan} terlepas: itulah \omterm{def:g5:water-shapes-the-land:weathering}{pelapukan}. Kepingan itu
jatuh ke \omterm{def:g1:landscapes-around-us:water}{sungai} kecil di bawahnya. Di aliran \omterm{def:g5:water-shapes-the-land:source-mouth}{hulu} airnya deras dan curam:
air menggulingkan kepingan itu (\omterm{def:g5:water-shapes-the-land:erosion}{erosi}), dan karena berbenturan dengan batu
lain, kepingan itu pecah menjadi \omterm{def:g3:stones-pebbles-sand:pebble}{kerakal} lalu menjadi butir-butir \omterm{def:g3:stones-pebbles-sand:sand}{pasir}. Di
aliran tengah \omterm{def:g1:landscapes-around-us:water}{sungai} melambat; butir itu beristirahat sebentar di pantai
berpasir di sisi dalam tikungan \omterm{def:g5:water-shapes-the-land:meander}{meander}, sampai \omterm{def:g5:water-shapes-the-land:flood}{banjir} mengangkatnya lagi.
Di aliran hilir \omterm{def:g1:landscapes-around-us:water}{sungai} itu lambat dan lebar; butir itu terbawa air sampai ke
muara. Di sana airnya hampir berhenti, dan butir itu mengendap bersama
\omterm{def:g3:stones-pebbles-sand:sand}{lumpur} dan \omterm{def:g3:stones-pebbles-sand:sand}{pasir} \omterm{def:g5:water-shapes-the-land:delta}{delta}: itulah \omterm{def:g5:water-shapes-the-land:deposition}{pengendapan}.
\end{solution}
